\chapter{Родитель космологической постоянной из сквозного потока}
\label{ch:v10-cosmological-throughflow-parent}

\section{Диссипативный источник кривизны}

Предыдущий гейт показал, что независимо выведенная космологическая постоянная
зафиксировала бы проводимость, но сама кривизна роста пока вычислялась из
$\kappa$. Проверим более физическую гипотезу: космологическая кривизна
является откликом на необратимый сквозной поток.

Для канонического хопфовского цикла уже выведены
\begin{equation}
 F_\circlearrowright=3\log2,
 \qquad
 \sigma_\circlearrowright=\kappa\log2.
 \label{eq:v10-lambda-flow-inputs}
\end{equation}
Единственная комбинация этих величин, имеющая размерность кривизны без
введения дополнительной длины, имеет вид
\begin{equation}
 \Lambda_{\rm flow}
 =\frac{3\sigma_\circlearrowright^2}
 {c^2F_\circlearrowright^2}.
 \label{eq:v10-lambda-flow-response}
\end{equation}
После подстановки точных данных цикла получаем
\begin{equation}
 \Lambda_{\rm flow}=\frac{\kappa^2}{3c^2}.
 \label{eq:v10-lambda-flow-conductance}
\end{equation}
Таким образом, искомая космологическая формула допускает не только
кинематическое, но и термодинамическое чтение: кривизна пропорциональна
квадрату производства энтропии, нормированному силой цикла.

\section{Условный положительный родитель}

Введём безразмерные переменные
\begin{equation}
 s:=\frac{\sigma_\circlearrowright}{\kappa\log2},
 \qquad
 u:=\frac{\Lambda c^2F_\circlearrowright^2}
 {3\sigma_\circlearrowright^2}.
\end{equation}
Минимальный родитель отклика равен
\begin{equation}
 \mathcal P_{\Lambda\mid\rm flow}(s,u)
 =\frac12(s-1)^2+\frac12(u-s)^2.
 \label{eq:v10-lambda-flow-parent}
\end{equation}
Он ограничен снизу и имеет единственный нуль $s=u=1$. Его гессиан
\begin{equation}
 \operatorname{Hess}\mathcal P_{\Lambda\mid\rm flow}
 =\begin{pmatrix}2&-1\\-1&1\end{pmatrix},
 \qquad \rank=2,
 \qquad \det=1
 \label{eq:v10-lambda-flow-hessian}
\end{equation}
положительно определён. При выключении потока
$\kappa=0$ получаем $\sigma_\circlearrowright=0$ и
$\Lambda_{\rm flow}=0$: поддерживаемая кривизна схлопывается вместе с
процессом прохождения.

\section{Граница происхождения}

Аффинность и производство энтропии уже выведены предыдущими гейтами.
Однако член, отождествляющий нормированную кривизну $u$ с нормированным
потоком $s$, отсутствует в унаследованном родительском функционале. Поэтому
\eqref{eq:v10-lambda-flow-parent} является допустимой условной архитектурой,
но пока новым законом отклика.

Кроме того, в логарифмических координатах $(\kappa,\sigma,\Lambda)$ карта
\begin{equation}
 C_{\Lambda\mid\rm flow}=
 \begin{pmatrix}-1&1&0\\0&-2&1\end{pmatrix}
 \label{eq:v10-lambda-flow-scale-map}
\end{equation}
имеет ранг $2$, ядро размерности $1$ и уничтожает вектор $(1,1,2)^T$.
Преобразование
\begin{equation}
 (\kappa,\sigma,\Lambda)\longmapsto
 (a\kappa,a\sigma,a^2\Lambda)
\end{equation}
не меняет относительный родитель. Динамический отклик объясняет наличие и
исчезновение кривизны, но не её абсолютную величину.

\begin{theorem}[Условная кривизна сквозного потока]
Положительный родитель $\mathcal P_{\Lambda\mid\rm flow}$ выбирает
$\Lambda_{\rm flow}=3\sigma_\circlearrowright^2/
(c^2F_\circlearrowright^2)=\kappa^2/(3c^2)$. Кривизна положительна при
ненулевом потоке и исчезает вместе с ним.
\end{theorem}

\begin{nogocriterion}[Закон отклика не должен вводиться вручную]
Термодинамическая совместимость формулы не доказывает её происхождение.
Пока член $u-s$ не получен из эйнштейновского, спектрального или иного уже
допущенного сектора общего родителя, он является новой условной связью.
Даже после его допуска сохраняется одна абсолютная масштабная мода.
\end{nogocriterion}

\begin{protocol}[Статус родителя кривизны потока]\begin{enumerate}
\item аффинность и энтропия цикла: \textbf{$2/2$};
\item положительный условный родитель: \textbf{построен};
\item условное происхождение отклика: \textbf{$5/5$};
\item унаследованные источники родителя: \textbf{$2/3$};
\item происхождение связи кривизна--поток: \textbf{$0/1$};
\item физическое происхождение $\Lambda$: \textbf{$0/1$};
\item абсолютный масштаб: \textbf{$0/1$};
\item следующий гейт: \textbf{происхождение связи из эйнштейновского
геометрического отклика}.
\end{enumerate}\end{protocol}

Машинный сертификат сохранён в
\path{s2t_v10_cell_birth_four_volume_cosmological_constant_throughflow_parent_origin_gate_results.json}.